\documentclass{article} % For LaTeX2e
\usepackage{iclr2027_conference,times}

% Optional math commands from https://github.com/goodfeli/dlbook_notation.
\input{math_commands.tex}

\usepackage{hyperref}
\usepackage{url}
\usepackage{booktabs}
\usepackage{multirow}
\usepackage{array}
\usepackage{xcolor}
\usepackage{graphicx}
\usepackage{amsmath,amssymb}
\usepackage{microtype}
\usepackage{enumitem}

% ---------------------------------------------------------------------------
% Draft conventions.  \ph{...} marks a number or phrase still to be filled from
% the full evaluation; \pilot marks a value measured on the five-episode
% long-tier pilot.  Both compile away to plain text once \finaldraft is set.
% ---------------------------------------------------------------------------
\definecolor{phbg}{HTML}{FFF1CC}
\definecolor{phfg}{HTML}{8A5A00}
\definecolor{pilotfg}{HTML}{1D6B3A}
\newif\iffinaldraft
%\finaldrafttrue
\iffinaldraft
  \newcommand{\ph}[1]{#1}
  \newcommand{\pilot}{}
\else
  \newcommand{\ph}[1]{\colorbox{phbg}{\textcolor{phfg}{\small\texttt{#1}}}}
  \newcommand{\pilot}{\textcolor{pilotfg}{\textsc{\scriptsize pilot}}}
\fi
\newcommand{\revoke}{\textsc{Revoke}}
\newcommand{\Ct}{\ensuremath{C_t}}
\newcommand{\Rt}{\ensuremath{R_t}}
\newcommand{\Et}{\ensuremath{E_t}}
\newcommand{\Dt}{\ensuremath{D_t}}
\newcommand{\pred}[1]{\texttt{#1}}

\title{\revoke: Does an Agent's Experience Expire\\When the Rules Do?}

% Authors must not appear in the submitted version. They should be hidden
% as long as the \iclrfinalcopy macro remains commented out below.
% Non-anonymous submissions will be rejected without review.

\author{Wentao Gao \\
Department of Computer Science and Engineering\\
The University of Texas at Arlington\\
Arlington, TX 76019, USA \\
\texttt{wentao.gao@uta.edu} \\
\And
\ph{co-authors} \\
\ph{affiliation} \\
\ph{address} \\
\texttt{\ph{email}}
}

\newcommand{\fix}{\marginpar{FIX}}
\newcommand{\new}{\marginpar{NEW}}

%\iclrfinalcopy % Uncomment for camera-ready version, but NOT for submission.
\begin{document}

\maketitle

\begin{abstract}
As LLM agents are deployed over long horizons, their memory systems have split into two families: declarative middleware that records facts about users and the world, and self-evolving experiential memory that distils insights, workflows and skills from the agent's own trajectories. Existing benchmarks probe the first family's recall fidelity and score answers; none asks what happens to either family when the rules of the world change. We present \revoke, a benchmark in which agents execute tool-use tasks in a programmatically verifiable environment whose constraints evolve across sessions through six event types: addition, supersession, conditionalisation, conflict, support and partial retraction. Grading is behavioural and deliberately asymmetric: an action that contradicts the current constraint closure is a violation, and everything else the agent does is unconstrained. Ground truth is a two-layer construction, a persistent base of defeasible rules with priorities and a per-step derivation layer that computes the grounded extension of a defeat graph followed by a stratified closure, giving unique, polynomial-time, deletion-order-independent grading with no LLM judge. Scenarios are generated by traversing a DAG of reusable event-subgraph motifs and laid onto realistic records of 300k to 1M tokens---a ward log, a chat export, a minute book---so that the constraint history is rich, verifiable and produced at scale. We name the failure mode under test \emph{procedural memory invalidation}: experience consolidated under an old world state becomes a violation-inducing rule after the state evolves, so that remembering more can mean violating more. In a five-episode pilot on the long tier, models with the whole transcript in context violate a current rule on 34--74\% of tasks in the flash tier and 17--21\% in the pro tier; a frontier model violates on 10\% once the licensed option can no longer be found by looking for what was never restricted. Across \ph{n} systems, we find \ph{headline finding}.
\end{abstract}

\section{Introduction}

Agent memory research advances along two largely disjoint tracks. The first treats memory as a declarative knowledge store: extract profiles, preferences and events from conversation, write them to external storage, retrieve on demand. Middleware such as mem0, MemOS and Zep belongs here~\citep{chhikara2025mem0,li2025memos,rasmussen2025zep}, evaluated by long-horizon question answering in the style of LoCoMo and LongMemEval~\citep{maharana2024locomo,wu2024longmemeval}: did the system remember correctly? The second treats memory as procedural experience. ExpeL distils cross-task insights, Agent Workflow Memory induces reusable workflows, Voyager and SkillWeaver synthesise skills~\citep{zhao2024expel,wang2024awm,wang2023voyager,zheng2025skillweaver}, and MemEvolve lets the memory architecture itself evolve~\citep{memevolve}. These systems are evaluated by task success on GAIA- and WebArena-style suites~\citep{mialon2023gaia,zhou2023webarena}: did the agent get better at doing? Both tracks sidestep the same question. \emph{When the rules of the world change, does what memory holds remain valid?}

For the declarative family the question has only just been raised. StateMemBench~\citep{statemembench} defines state tracking across sessions, where facts and decisions are revised and answers must reflect the current rather than a superseded state. This is an important step, but its grading operates at the answer level. Real agent failures occur at the level of behaviour: booking a vendor that has since been banned, calling an API that was deprecated, administering a drug that acquired a contraindication. A correct answer does not imply compliant behaviour, and behavioural grading naturally accommodates open-ended, multi-solution execution without a fixed answer space.

For the experiential family the question has never been examined. Procedural experience carries an implicit timestamp: every distilled insight and consolidated workflow was validated under some particular world state. When constraints evolve, that experience does not expire on its own. It continues to be retrieved and reused, and because it remains effective on the old distribution it can survive selection in success-driven evolution. We call this failure mode \textbf{procedural memory invalidation}. A system that only accumulates and never invalidates will exhibit the counterintuitive pattern of \emph{more experience, more violations}.

\revoke\ evaluates both families under one behavioural criterion. Our contributions:
\begin{enumerate}[leftmargin=*]
\item We formalise procedural memory invalidation and move evolving-state evaluation from the answer level to the behaviour level, grading traces rather than responses (\S\ref{sec:formal}).
\item We give a two-layer constraint construction---a grounded extension over a defeat graph followed by stratified Datalog closure---that guarantees unique, polynomial-time, deletion-order-independent grading and sidesteps the over-deletion pitfall of incremental view maintenance.
\item We build the benchmark as a \emph{graph traversal}: reusable event-subgraph motifs are composed by randomised topological linearisation into long, verifiable constraint histories, with ground truth always recomputed by the solver rather than taken from the template's intent (\S\ref{sec:construction}). Three tiers separate memory-truncation failures, reasoning failures under noise and speaker hierarchy, and failures of memory maintenance over 300k--1M tokens of realistic transcript.
\item We introduce metrics beyond pass/fail: a difficulty weight derived from the state history, an exam score that penalises acting wrongly below not acting, and per-event violation, adaptation lag, interference and staleness attribution, all computable on the graph without annotation (\S\ref{sec:protocol}). We report the first head-to-head comparison of declarative and experiential memory systems under this criterion (\S\ref{sec:experiments}).
\end{enumerate}

\section{Formalisation: a two-layer graph construction}
\label{sec:formal}

We represent constraint state as a graph, but split the intuitive picture---rules as nodes, relations as edges, contradiction as deletion---into two layers. Deletion happens only in a derivation layer recomputed at every timestep; the persistent layer never loses information on conflict. The split gives ``delete the contradicted node'' an exact semantics: it is precisely the iterated computation of the grounded extension of an argumentation framework in the sense of \citet{dung1995acceptability}, unique, polynomial-time, and independent of deletion order.

\paragraph{Persistent layer.}
Ground truth at time $t$ is a rule base $\Rt = \{(r_i, \pi_i)\}$ of defeasible rules with explicit integer priorities over a decidable Datalog fragment: heads are literals with classical negation, bodies are conjunctions of positive literals, variables are sorted. Six events drive evolution. \textsc{Add} inserts a rule. \textsc{Supersede} replaces a rule's head or body. \textsc{Condition} rewrites a body between unconditional and conditional forms. \textsc{Conflict} inserts a rule with a contradictory head, resolved by priority rather than deletion. \textsc{Support} inserts nothing and increments the priority of its target, folding support into priorities and avoiding bipolar frameworks. \textsc{Retract} withdraws named ground instances of a rule. Except for \textsc{Supersede} and \textsc{Retract} the layer is append-only, which is what makes staleness attribution well defined.

\paragraph{Derivation layer.}
Recomputed from scratch at every timestep and never stored (Table~\ref{tab:derivation}). Nodes are ground \emph{instances}, not rules, so a rule defeated in one context stays valid in another. Because the closure is recomputed over the whole surviving set rather than by deleting a node and its outgoing edges from the old closure, any conclusion with an alternative derivation is automatically re-derived: the situation Delete-and-Rederive exists to handle in incremental Datalog~\citep{gupta1993maintaining}, which we sidestep by per-step recomputation.

\begin{table}[t]
\centering\small
\begin{tabular}{@{}p{2.3cm}p{2.2cm}p{8.4cm}@{}}
\toprule
Step & Output & \\
\midrule
1 Ground & $\Rt \to N_t$ & Instantiate every live rule over the sorted Herbrand universe. \\
2 Defeat graph & $n \to m$ & An edge whenever two instances have complementary heads and $\pi(n) > \pi(m)$. Priorities are strict, so the relation is acyclic. \\
3 Grounded ext. & $\Et$, $\Dt$ & Iterated deletion to fixpoint: a node is in once all its attackers are out, out once any attacker is in. Survivors $\Et$, deleted $\Dt$. \\
4 Closure & $\Ct$ & Stratified forward chaining over $\Et$ alone, recomputed over the whole surviving set. \\
\bottomrule
\end{tabular}
\caption{The derivation layer.}
\label{tab:derivation}
\end{table}

\paragraph{Grading and properties.}
An action $a$ asserts a set of ground literals. It is a \textbf{violation} iff $\Ct \cup \mathrm{assert}(a)$ derives a contradiction, where the asserted literals are pushed through the surviving instances, so a violation can surface several derivation steps away from the action itself. Three properties hold throughout evolution. (i) $\Et$ is unique and polynomial-time computable. (ii) Removing nodes or edges cannot create negative cycles, so every subgraph remains stratifiable; the semantics is well defined at every timestep. This is well-definedness, not conclusion preservation: negation makes the logic non-monotone and deletion can flip conclusions, which is exactly the phenomenon under test. (iii) The permissive default: if the closure says nothing about an entity, acting on it is not a violation.

\paragraph{Worked example.}
Procurement, after session 5: a fact $f_1$: \pred{low\_budget}; $r_1$: \pred{ok(A)} $\Leftarrow \top$ at $\pi{=}1$ (\textsc{Add}, session 2, the team lead's default vendor); $r_2$: $\neg$\pred{ok(A)} $\Leftarrow$ \pred{low\_budget} at $\pi{=}2$ (\textsc{Conflict}, session 5, compliance). The derivation has the edge $n_2 \to n_1$, so $\Et = \{f_1, n_2\}$, $\Dt = \{n_1\}$, $\Ct = \{\pred{low\_budget}, \neg\pred{ok(A)}\}$: ordering $A$ asserts \pred{ok(A)}, a contradiction, a violation; ordering $B$ meets a silent closure and passes. In session 9 a \textsc{Support} event---an executive directive endorses $A$---raises $\pi(r_1)$ from 1 to 3 without changing any rule text. The edge reverses, $\Et = \{f_1, n_1\}$, and ordering $A$ is compliant again. A memory system that froze the session-5 conclusion ``$A$ is banned'' as a fact is now stale, while a system storing rules plus priorities adapts for free. The engine reproduces this example as an executable test.

\section{Benchmark construction}
\label{sec:construction}

\subsection{The generator is a graph traversal}
Constraint evolution is not hand-written per scenario. A \textbf{motif} is a small, domain-agnostic event subgraph over a permission predicate \pred{allow/1}, a set of 0-ary context predicates and a group predicate \pred{grp/2}. Each motif emits an ordered chain of \textbf{beats}; a beat is one atomic step of evolution, optionally followed by a \textbf{probe}, a task the agent must perform at that moment. Chains are internally ordered and mutually independent, so the beats of all motifs in a scenario form a DAG whose linear extensions are the admissible timelines. The generator draws a randomised topological linearisation, weighting each chain by how many beats it has left, and lays the result onto a session grid with filler turns, distractor options and echo probes.

\textbf{Ground truth is never taken from a motif's design intent.} Once the timeline exists, every probe is re-scored against the closure the solver actually produces at that session. An interleaving that changes the intended semantics is therefore caught by the verifier rather than silently mislabelled. Every motif is designed so that the option set offered at consecutive probes stays constant while the compliant subset changes: an agent replaying experience from an earlier probe violates.

\begin{table}[t]
\centering\small
\begin{tabular}{@{}lp{4.6cm}p{6.2cm}@{}}
\toprule
Motif & Shape & What a stale memory gets wrong \\
\midrule
\multicolumn{3}{@{}l}{\emph{Easy tier}} \\
conflict\_flip & \textsc{Add} $\to$ \textsc{Conflict} $\to$ \textsc{Support} & a banned entity becomes permitted again while its old alternative is withdrawn \\
supersede\_alt & \textsc{Add} $\times 2$ $\to$ \textsc{Supersede} & the consolidated path is retired; the task survives via a second derivation \\
conditionalize & \textsc{Add} $\to$ \textsc{Conflict} $\to$ \textsc{Condition} $\to$ context fires & a blanket ban is narrowed, then re-triggered by context \\
retract\_partial & \textsc{Add} $\to$ group ban $\to$ \textsc{Retract} & a member leaves the blanket ban exactly as the old carve-out closes \\
conflict\_chain & three-way conflict, transitive priority & the middle rule's verdict is overturned by a higher one \\
support\_threshold & losing objection is \textsc{Support}ed past the permission & no rule text changes, only a priority \\
condition\_widen & conditional ban $\to$ unconditional & the qualifier the agent relied on disappears \\
canary & no applicable constraint & calibrates the false-positive rate \\
\midrule
\multicolumn{3}{@{}l}{\emph{Hard and long tiers (add speaker ranks, near-miss noise, long arcs)}} \\
ctx\_flipflop & conditional ban; context on/off/on; then widened & a cached ``X is banned'' is wrong half the time \\
alias & ``X follows Y''; Y banned, restored, banned; link broken & transitive effects, then their absence \\
group\_dynamics & conditional group ban; members join and leave & membership stated dozens of sessions earlier \\
hierarchy & rank-3 ban, later rank-2 ``default'', lift, rank-2 ban & recency is not authority \\
stale\_reminder & rank-3 ban, then unranked ``reminders'' that X is the default & statements that look like rules but are not \\
reinstate\_arc & retire $\to$ reverse $\to$ conditionalise $\to$ context on/off & a six-step arc on one entity \\
proposal\_noise & proposals and hearsay that never become rules & treating a proposal as a ban \\
threshold & ban above a numeric limit; the limit moves & a threshold compared against its old value \\
derived & pairwise incompatibility; the partner becomes active, then inactive & a prohibition that is never stated, only derived \\
conjunctive & ban only while two conditions hold together & one condition treated as sufficient \\
certification & a list regime switched on and off; the list amended & off-list use treated as merely unlicensed \\
\bottomrule
\end{tabular}
\caption{The motif library. Each is stated only in the engine's vocabulary, so every world is graded by the same code. The certification regime is an episode-level rank-3 rule: while it is on, an option not on the list is a violation, which is what stops ``avoid the banned ones and pick anything else'' from completing a task.}
\label{tab:motifs}
\end{table}

\subsection{Worlds and tiers}
Five seed domains---platform API governance, procurement compliance, clinical medication constraints on a fully synthetic formulary, household automation preferences, financial operation permissions---each supply 24 entities (48 in the hard tier), six context predicates, three groups, a mock tool surface with one graded \emph{act} tool and read-only tools that deliberately return nothing substantive, and natural-language templates for every event type. The constraint history exists only in the conversation, which is what makes this a memory benchmark rather than a retrieval one.

The \textbf{easy tier} uses 10--40 sessions, a single team voice and obviously irrelevant filler. The \textbf{hard tier} adds five pressures that keep every verdict derivable from the transcript: 60--140 sessions with several motifs in flight; a declared \emph{speaker hierarchy} in which a higher-ranked ruling beats a lower-ranked one regardless of order, so recency stops being a valid heuristic; \emph{near-miss noise}---proposals, hearsay, questions, cross-team anecdotes and stale ``reminders'' from speakers the hierarchy declares to have no authority, aimed at the current state; probes asked by unranked speakers who nudge towards a currently forbidden option; and terse, referential updates (``our ruling from session 16 is withdrawn'').

\subsection{The long tier}
\label{sec:long}
The long tier keeps the hard tier's logic and changes the medium. Each episode is a single continuous record of 420--650 sessions, rendered as the document such a history would actually live in and padded to a target of 300k--1M tokens. A world is a JSON \emph{corpus} that supplies vocabulary and surface only: 96 fictional entities, six conditions, six groups, phrasings for every event kind, people by rank with their declared authority terms, eight kinds of near-miss noise, a numeric threshold, a certification list, and a slot-template filler in the world's own register. The engine, the events and the grader are untouched. Six worlds are built: a ward round and nursing handover log for one long-stay patient; a chat export of a \texttt{\#platform-release} channel; design-crit notes merged with a design-system decision log; the group chat of an older adult's care circle; a live-service game studio's weekly narrative review; and a company minute book produced by an AI notetaker.

Two layouts are used. In a \emph{sectioned} record the rule events go under the decisions heading attributed to their speaker, chatter goes under discussion---including lines phrased exactly like recorded decisions by people with no authority---and the task goes under action items; the section header is a cue and the phrasing a counter-cue, and only the speaker settles it. In a \emph{flat} record rules, chatter and filler share one stream with a speaker on every line and nothing marking which is which. Filler is drawn from slot templates whose product is on the order of $10^9$ distinct lines; two invariants, enforced by a validator and again by the padder, make it safe for grading: filler never names a governed entity, and never uses the rule register. A third rule is \textbf{rank neutrality}: the engine decides who speaks a rule template, often a rank-1 person, so no template may assert an authority (``Council ruling: \ldots'') that the grader will not grant its speaker.

A released episode carries 10--15 of the 180--270 tasks the generator emits: a fifth from the first third of the record, three tenths from the middle, half from the last third, preferring stale-memory traps to about seven in ten, spread over event types, never two in one session. The other task turns are removed, so the record reads as months of conversation with a handful of moments where the agent is asked to act; the state history behind each kept task is unchanged, so its labels stay valid. Every generated task is retained in the released item as a \emph{task pool} with its labels and rendered text, so a different selection can be asked of the same episode without regenerating it. Two rules on option sets came out of the pilot (\S\ref{sec:pilot}): the always-approved \emph{anchor} entities that keep every task completable carry their own arcs of prohibition and reinstatement, and the selector avoids any task whose licensed options were all approved long ago and never restricted since.

\subsection{Two grading gates}
Violations are graded permissively. That alone is exploitable twice over: \emph{never act}, or \emph{always act on something nobody ever mentioned}. Both are violation-free and both are useless. A task therefore counts as \emph{done} only if the agent acted on an option the closure positively licenses ($\pred{allow}(e) \in \Ct$); it counts as \emph{violated} if the action contradicts the closure; everything else is an \emph{abstention}. Every probe offers 3--6 options mixing entities the transcript permits, prohibits and never mentions, in shuffled order. The scripted \emph{refuse} policy makes the point: it is violation-free and completes nothing.

\subsection{Solver-side acceptance}
\label{sec:acceptance}
A scenario is emitted only if it passes seven checks, and failures are rejected rather than repaired: (1) the closure is convergent and consistent at every session; (2) no probe is must-violate, and every probe has a licensed option; (3) the over-deletion guard: a feasibility predicate with one derivation per permitted entity is present exactly when some permission survives; (4) at least one stale-memory trap; (5) every probe is rendered into the transcript; (6) self-containedness: every violating option was named in an earlier turn; (7) recoverability: whenever a surviving instance defeats a deleted one, the transcript licenses that outcome by recency, an explicit override, an advisory marker, or the declared speaker hierarchy. Acceptance rates are \ph{xx}\% (easy), \ph{xx}\% (hard), and about one seed in five on the long tier once anchors can be prohibited (2.4--4.8 seeds rejected per accepted episode in the release build).

\begin{table}[t]
\centering\small
\begin{tabular}{@{}lrrrrrr@{}}
\toprule
 & Episodes & Sessions/ep. & Tokens/ep. & Graded decisions & Stale traps & Recall span $>200$ \\
\midrule
Easy tier & \ph{1000} & \ph{36.7} & \ph{n} & \ph{17,702} & \ph{6,672} & -- \\
Hard tier & \ph{n} & \ph{$\sim$130} & \ph{n} & \ph{n} & \ph{n} & \ph{n} \\
Long tier, pilot~\pilot & 5 & 446 & 368k--1018k & 50 & 36 & 40 \\
Long tier, release & 100 & 542 & 344k--934k & 1202 & 848 & 971 \\
\bottomrule
\end{tabular}
\caption{Dataset composition. The pilot episodes, in order of length: ward log 368k (sectioned, 439 sessions), release channel 500k (flat, 444), design log 665k (sectioned, 443), care-circle chat 850k (flat, 458), narrative review 1018k (sectioned, 445). Prompt tokens measured at the API are within 3\% of the character-based estimate. The release is 100 episodes of 10--15 tasks over five length bands (300k, 440k, 580k, 720k, 830k targets; 6--10 motif cycles; 428--671 sessions), 16--17 per world, built in one run with no failed seed; 19 of the 1202 released tasks are of the never-restricted kind. Failures within an episode are serially correlated, so the episode is the unit of analysis (\S\ref{sec:limitations}).}
\label{tab:dataset}
\end{table}

\section{Evaluation protocol}
\label{sec:protocol}

An episode is replayed session by session. Update, noise and filler turns are delivered as messages; at a probe the agent must call the act tool with one of the offered options. Read tools are free and ungraded. The grader rebuilds ground truth from the item's event log with the same engine that generated it and never trusts cached labels. Only the blind view of an item is ever exposed to an agent: transcript text, tool surface and option sets, with turn kinds stripped so a scaffold cannot filter updates from noise by construction. The harness is plain for every system: the agent is asked to act and to say in a sentence what it did.

\paragraph{Memory conditions.}
\emph{full} keeps the whole transcript in context, the long-context upper bound. \emph{compact-$K$} is one agent running the whole episode continuously with a context budget of $K$ tokens: as each session arrives and the budget overflows, the oldest material leaves the context for good and the agent rewrites it into running notes it must later act on, with no access to the original text. The notes are bounded to half the budget, a fraction of the budget is kept verbatim, and compaction happens at session granularity rather than only when a task comes up. \emph{truncate} drops the oldest material without notes; it is an ablation, not a baseline. \emph{middleware} conditions route the transcript through mem0, MemOS, Zep and OmniMemory; \emph{experiential} conditions let ExpeL, AWM, Dynamic Cheatsheet~\citep{suzgun2025dynamic} and Voyager-style skill libraries consolidate across episodes of the same world~\ph{(planned)}.

\paragraph{Grading a trace without a judge.}
The trace is never read. One structured action is taken out of it and one decidable question is asked about that action, in two steps. First, the act tool's argument is resolved to an entity: an exact name or id, ignoring case and a leading article, or a short single-line phrase naming exactly one of the 96 entities. This is a lookup against a known list, not language understanding, and it is possible only because the option set is closed. Anything else resolves to nothing and counts as an abstention, never a violation. Second, the grader rebuilds $\Rt$ from the event log up to that session, recomputes $\Ct$, and asserts $\pred{allow}(e)$ into it: a derivable contradiction is a \emph{violation}, $\pred{allow}(e) \in \Ct$ is a \emph{completion}, and neither is an \emph{abstention}. All three are queries over a Datalog closure, so nothing in the pipeline depends on what the agent wrote, and no LLM judge appears anywhere.

The strictness of the first step was forced by a measurement failure. An earlier version substring-matched the argument and took the longest hit; compacting agents often refuse by writing their reasoning into the argument slot, enumerating every option before rejecting it, and the matcher scored those refusals as actions --- usually violating ones, since the longest name in an enumeration is as likely to be banned as not. This affected 53\% of the steps in one condition and inflated its violation rate from 0.048 to 0.462, inverting two findings. Unresolved arguments are now counted in every report so the artefact cannot recur silently.

\paragraph{What can be said about the agent's reasoning.}
Whether an action rested on a premise that has since been overturned is answered structurally, not by reading the rationale: the grader records, for each violation, the rule instances that make the choice a violation, and whether the permission the action needed is an instance in the \emph{deleted} set $\Dt$ --- a licence that existed once and has been defeated. Together with the stale-trap label and the recall span, this covers most of what acting on a stale premise means, with no judge. Two things it does not cover. Requiring the agent to emit a status and a justifying session for every option would make the rationale checkable line by line against the closure, still without a judge, but it is a chain-of-thought scaffold that almost no deployed agent receives and it makes frontier models near-perfect; we therefore fix a plain harness for the main results and keep structured elicitation as an ablation. Whether a free-text rationale faithfully reports the agent's grounds is outside the closure entirely; any such analysis would have to be secondary, on a sample, with inter-judge agreement, and never feeding the headline metrics, since a judge faces the same long-context retrieval task the benchmark exists to measure and its errors would correlate with the failures under study.

\paragraph{Metrics.}
\emph{Violation rate}: share of tasks where the action contradicts the current closure, reported absolute and difficulty-weighted. \emph{Completion} and \emph{abstention}: share of tasks where a licensed option was chosen, and share where nothing licensed and nothing violating was chosen; violation alone hides the second. \emph{Exam score}: each task scores $+1$ (licensed), $0$ (abstained) or $-1$ (violated), weighted and averaged into $[-1, 1]$; negative means worse than never acting. \emph{Difficulty weight}, from ground truth alone:
\begin{equation}
w = \max\!\big(0.3,\; 1 + 1.0\,[\mathrm{span} > 200] + 0.8\,[\mathrm{trap}] + 0.4\,(n_{\mathrm{viol}} - 1)^{+} - 0.3\,(n_{\mathrm{lic}} - 1)^{+}\big),
\end{equation}
where span is the number of sessions back to the oldest status change among the offered options that still governs the verdict. Violation is further decomposed by event type and by recall span (near $\le 96$, mid $\le 200$, far $> 200$ sessions). \emph{Stale-trap rate}, \emph{adaptation lag}, \emph{update interference} and \emph{staleness attribution} (share of violations whose supporting permission instance lies in $\Dt$) follow the design; canary probes calibrate false violations and over-refusal. Position in the episode is deliberately not a difficulty feature: it correlates with violation about as strongly as recall span, but only because the generator places the harder event kinds later.

\section{Scripted reference policies}
\label{sec:scripted}

Five policies with no model calls establish the floor and ceiling. \emph{oracle} always picks a licensed option. \emph{stale} acts on the option that was compliant at the previous task and is a violation now, when there is one: procedural memory invalidation in its purest form. \emph{recency} picks the option named in the most recent rule. \emph{random} is uniform over the offered options. \emph{refuse} never acts.

\begin{table}[t]
\centering\small
\begin{tabular}{@{}lrrrrr@{}}
\toprule
Policy & Violation & Completion & Abstention & Exam & Trap violation \\
\midrule
oracle  & 0.00 & 1.00 & 0.00 & $+1.00$ & 0.00 \\
stale   & 0.79 & 0.21 & 0.00 & $-0.63$ & 0.96 \\
recency & 0.84 & 0.16 & 0.00 & $-0.69$ & 0.84 \\
random  & 0.63 & 0.23 & 0.14 & $-0.44$ & 0.62 \\
refuse  & 0.00 & 0.00 & 1.00 & $\phantom{+}0.00$ & 0.00 \\
\bottomrule
\end{tabular}
\caption{Scripted reference policies on the long tier, all 1202 released tasks of the 100-episode release. Random violates two tasks in three because option sets are tight: on average one licensed option among five. The stale policy's trap violation of 0.96 is the number the benchmark exists to produce.}
\label{tab:scripted}
\end{table}

\section{Experiments}
\label{sec:experiments}

\subsection{Systems}
Three baseline groups: long-context (full history in context), retrieval (RAG over past sessions), and memory systems, both declarative middleware (mem0, MemOS, Zep, OmniMemory) and self-evolving experiential systems (those unified in EvolveLab~\citep{evolvelab}, including ExpeL, AWM, Dynamic Cheatsheet and Voyager-style skill libraries). Backbones: \ph{model list}. Each system runs \ph{k} seeds; we report means and \ph{interval}.

\subsection{Main results}
\begin{table}[t]
\centering\small
\begin{tabular}{@{}lrrrrrrr@{}}
\toprule
System & Exam easy & Exam hard & Exam long & Viol.\ long & Stale trap & Lag & Attribution \\
\midrule
\multicolumn{8}{@{}l}{\emph{Long-context}} \\
\ph{backbone}~$\cdot$ full & \ph{x.xx} & \ph{x.xx} & \ph{x.xx} & \ph{x.xx} & \ph{x.xx} & \ph{x.x} & \ph{x.xx} \\
\ph{backbone}~$\cdot$ compact-8k & \ph{x.xx} & \ph{x.xx} & \ph{x.xx} & \ph{x.xx} & \ph{x.xx} & \ph{x.x} & \ph{x.xx} \\
\multicolumn{8}{@{}l}{\emph{Retrieval}} \\
RAG~$\cdot$ \ph{retriever} & \ph{x.xx} & \ph{x.xx} & \ph{x.xx} & \ph{x.xx} & \ph{x.xx} & \ph{x.x} & \ph{x.xx} \\
\multicolumn{8}{@{}l}{\emph{Declarative middleware}} \\
mem0 / MemOS / Zep & \ph{x.xx} & \ph{x.xx} & \ph{x.xx} & \ph{x.xx} & \ph{x.xx} & \ph{x.x} & \ph{x.xx} \\
\multicolumn{8}{@{}l}{\emph{Self-evolving experiential}} \\
ExpeL / AWM / Cheatsheet / skills & \ph{x.xx} & \ph{x.xx} & \ph{x.xx} & \ph{x.xx} & \ph{x.xx} & \ph{x.x} & \ph{x.xx} \\
\bottomrule
\end{tabular}
\caption{Main results (planned). Exam score on all three tiers; the remaining columns on the long tier.}
\label{tab:main}
\end{table}

\paragraph{Analyses (planned).} Per-event decomposition: \ph{which event types drive violations for each family; expected ordering \textsc{Support} and \textsc{Retract} hardest for conclusion-caching systems}. More experience, more violations: for experiential systems, violation rate against the number of consolidated episodes in the same world, \ph{slope and sign per system}. Adaptation lag and interference: \ph{distributions per family}. Staleness attribution: for systems with inspectable stores, \ph{attribution rate per system; share of stale entries written as conclusions rather than rules}.

\subsection{The long-tier pilot}
\label{sec:pilot}
Five episodes, one per world except the minute book, one seed each, fifty tasks~\pilot. Every model was run with the whole transcript in context, the condition most favourable to it; one flash-tier backbone was also run under two compaction budgets. Every violation reported below was checked by hand against the events that make the chosen option a violation, and every tool call was checked for the refusal-in-the-argument pattern (\S\ref{sec:protocol}): 1 of 147 flash-tier calls, none elsewhere.

\begin{table}[t]
\centering\small
\begin{tabular}{@{}lrrrrrrr@{}}
\toprule
Backbone, full context & Tasks & Violation & Completion & Abstention & Exam & Far-span viol. & Cost \\
\midrule
\multicolumn{8}{@{}l}{\emph{v2 episodes (anchors with arcs; released construction)}} \\
gpt-5.5 (368k and 500k episodes) & 20 & 0.10 & 0.90 & 0.00 & $+0.81$ & 0.12 & \$11.8 \\
deepseek-v4-flash & 50 & 0.34 & 0.44 & 0.22 & $+0.07$ & 0.40 & \$1.0 \\
glm-5.3-flash & 50 & 0.40 & 0.38 & 0.22 & $-0.04$ & 0.42 & \$1.3 \\
qwen3.7-flash & 47 & 0.74 & 0.19 & 0.06 & $-0.57$ & 0.78 & \$3.2 \\
\midrule
\multicolumn{8}{@{}l}{\emph{v1 episodes (constant anchors; superseded, see text)}} \\
gpt-5.5 (368k, 500k, 849k episodes) & 30 & 0.00 & 1.00 & 0.00 & $+1.00$ & 0.00 & \$29.9 \\
deepseek-v4-pro & 47 & 0.17 & 0.60 & 0.23 & $+0.41$ & 0.19 & \$4.2 \\
glm-5.3 & 48 & 0.21 & 0.62 & 0.17 & $+0.38$ & 0.21 & \$7.3 \\
deepseek-v4-flash & 47 & 0.36 & 0.40 & 0.23 & $+0.03$ & 0.43 & \$0.9 \\
glm-5.3-flash & 49 & 0.43 & 0.39 & 0.18 & $-0.04$ & 0.49 & \$1.7 \\
qwen3.7-flash & 47 & 0.66 & 0.23 & 0.11 & $-0.50$ & 0.76 & \$2.7 \\
\bottomrule
\end{tabular}
\caption{Full-context results on the pilot~\pilot. Task counts below 50 are the last probes of the 1018k episode, which did not fit a 1,048,576-token window with output headroom (nor qwen's 1,000,000 window). Costs at list price.}
\label{tab:pilot-full}
\end{table}

\paragraph{Three tiers of backbone, one ordering.} With the whole record in context and nothing to remember, flash-tier models violate a current rule on a third to three quarters of tasks and abstain on another fifth; pro-tier models of the same vendors violate on a fifth; a frontier model on a tenth. The ordering is the same on both versions of the episodes and across worlds. The exam score adds what the violation rate hides: glm-5.3-flash's $-0.04$ says that with the entire transcript available it did slightly worse than an agent that never acts.

\paragraph{What goes wrong is a standing rule with an on/off state.} The certification regime---a rank-3 rule stated in session 6--8 that while a flag is on only listed options may be used, with the flag switched on and off every few dozen sessions and the list amended---is behind 16 of glm-5.3-flash's 20 violations, 16 of deepseek-v4-flash's 17 and 30 of qwen's 35, and one of gpt-5.5's two. The models read the rule, watch the list change, and then act as if the flag were off. The remaining violations are plain stale rulings: gpt-5.5's other one was an endpoint banned in session 50 and chosen in session 406. A rank-0 requester's nudge towards a violating option (``refunds-async should be fine for this, no?'') was followed into a violation on 10 of 28 nudged tasks by glm-5.3-flash.

\paragraph{Difficulty comes from the record's structure and the state's history, not from length.} The 1018k narrative review is not the hardest episode for any model; for glm-5.3-flash it is the easiest (violation 0.30, exam $+0.41$). The two flat chat exports and the design log are the hardest: a rule that arrives as one line among a hundred, from a speaker whose rank has to be looked up, with nothing marking it as a decision, is harder to keep than the same rule under a ``Decisions'' heading, whatever the total length.

\begin{table}[t]
\centering\small
\begin{tabular}{@{}lrrrrrrr@{}}
\toprule
glm-5.3-flash, memory condition & Viol. & Weighted & Compl. & Abst. & Exam & Compactions/ep. & Context at task \\
\midrule
full context & 0.40 & 0.42 & 0.38 & 0.22 & $-0.04$ & 0 & 368k--1018k \\
compact-2k, keep 55\%, notes $\le$1.2k & 0.26 & 0.28 & 0.16 & 0.58 & $-0.12$ & 417 & 2.1k \\
compact-8k, keep 93\%, notes $\le$4k & 0.28 & 0.29 & 0.40 & 0.32 & $+0.10$ & 399 & 10.2k \\
\bottomrule
\end{tabular}
\caption{One agent running each v2 episode continuously, the material leaving its context rewritten into bounded notes it must later act on~\pilot. Fifty tasks per row. The compact-2k notes overran their cap by about a third on 82\% of compactions even after a condensing pass; total context at a task still averaged the budget. The compact-8k row sits above its budget because retained text and notes are counted together.}
\label{tab:pilot-compact}
\end{table}

\paragraph{A maintained memory beats the whole record, mostly by declining.} Both compaction budgets lower the violation rate below full context ($0.40 \to 0.26$--$0.28$), reproducing the shorter tier's finding that frequent small compactions beat a long raw context for a flash-tier model. The 2k budget achieves it largely by abstaining: 58\% of tasks end with the agent unable to name an approved option, and the exam score is below full context. The 8k budget with a 93\% keep fraction---about 400 compactions of 560 tokens each per episode---is the only condition in the pilot on which this backbone scores above zero. Under compaction 12 and 4 of the 50 answers were refusals written into the argument slot; the grader scores them as abstentions.

\paragraph{Notes grow without bound unless bounded.} Left to itself the compacting model added two to three thousand characters of notes per session (3.4k, 5.9k, 8.8k over three consecutive sessions of the 1018k episode), so that within a few dozen sessions the notes alone exceeded a 2k-token budget and the condition's label meant nothing. The protocol therefore caps the notes at half the budget, tells the model the word limit and what to drop first, and runs one condensing pass on overrun. The unbounded runs were discarded.

\paragraph{The anchor lesson.} To keep every task completable, the generator gives each episode eight \emph{anchor} entities, approved by a rank-2 speaker in sessions 2--4 and certified throughout. In the first build of the pilot those anchors were never mentioned by a rule again, and in 48 of 50 tasks the only licensed option was one of them. A model with the whole record in context can find such an entity by absence, and gpt-5.5 did exactly that: 30 of 30 tasks correct, choosing the perpetual anchor of each option set. That measured retrieval of what was never restricted, not tracking of what changed. In the released construction each anchor carries one or two arcs---a rank-3 prohibition later reversed by the same level---the certification list may drop anchors, and the task selector avoids any task whose licensed options were all approved more than a hundred sessions earlier and never restricted since: 1 of 50 released pilot tasks is of that kind, against 48 before. The flash tier was unaffected by the change (0.36--0.66 before, 0.34--0.74 after); the frontier model went from 0 to 2 violations in 20 tasks, both on the flat release channel, both more than 350 sessions after the rule they broke.

\paragraph{What the pilot does not establish.} Five episodes are five clusters, and failures within an episode are serially correlated (\S\ref{sec:limitations}): the numbers above order the backbones and locate the failure modes, and are not confidence intervals. The frontier model has not been run under compaction, the condition closest to deployment, because its output pricing makes four hundred note-writing calls per episode cost more than the whole flash-tier grid. The pro-tier rows are from the first build only.

\section{Related work}
\paragraph{Long-horizon memory benchmarks.} LoCoMo and LongMemEval measure recall over long conversations with answer-level scoring~\citep{maharana2024locomo,wu2024longmemeval}; StateMemBench adds revised state across sessions, still QA-shaped~\citep{statemembench}. \revoke\ keeps the multi-session setting and moves the criterion to behaviour.
\paragraph{Memory systems.} Declarative middleware~\citep{chhikara2025mem0,li2025memos,rasmussen2025zep} and self-evolving experiential memory~\citep{zhao2024expel,wang2024awm,wang2023voyager,zheng2025skillweaver,suzgun2025dynamic,memevolve}, unified for comparison in EvolveLab~\citep{evolvelab}. To our knowledge no prior work quantifies invalidation of consolidated experience.
\paragraph{Defeasible reasoning and argumentation.} The grounded extension is Dung's~\citep{dung1995acceptability}; our two-layer construction instantiates it over ground rule instances with strict priorities and composes it with stratified Datalog. Delete-and-Rederive and related view-maintenance algorithms~\citep{gupta1993maintaining} motivate the recomputation choice.
\paragraph{Agent safety and constraint following.} Rule-following and tool-use safety suites grade compliance against static rule sets; \revoke\ grades against a rule set that evolves under the agent, which is where memory enters.

\section{Limitations}
\label{sec:limitations}
\begin{itemize}[leftmargin=*]
\item Constraint prose and filler are template-generated. The long tier's filler is drawn from a slot product large enough that verbatim repetition is rare (below 2\% of filler lines in the release build), but the register is uniform within a world; an LLM-assisted instantiation pass over the same event graph would raise surface realism without touching the ground truth.
\item The episode is the unit of analysis. On the shorter tier, failures under compaction came in runs of up to 55 consecutive tasks, and the probability of failing given the previous task failed was 0.92 against a base rate of 0.79, so a 180-task episode is worth roughly fourteen independent observations. The pilot's five episodes order systems; they do not bound them. The release is 100 episodes for that reason, and every released episode also carries its full task pool so that other selections can be drawn from the same history.
\item Probes offer an explicit option set, buying exact discriminativeness at the cost of some open-endedness.
\item Frontier backbones have been run only with the whole record in context; under compaction they would cost an order of magnitude more per episode at current output prices.
\item Echo probes repeat a task at consecutive sessions to make adaptation lag observable; they are excluded from the long tier's released tasks.
\item Update interference is negative for policies that replay their own history; the metric is read with its two arms.
\end{itemize}

\section{Conclusion}
\revoke\ asks a question neither memory family has been tested on: when the rules change, does what memory holds remain valid, and does the agent's behaviour show it? By grading traces against a constraint closure recomputed at every step, and by generating long, verifiable constraint histories as traversals of an event-subgraph DAG laid onto realistic million-token records, the benchmark isolates memory consistency from policy optimality and puts a number on procedural memory invalidation. The pilot already shows the shape of the answer: given the entire record, current flash-tier models break a standing rule on a third to three quarters of tasks, and the most common way to do it is to forget that a rule has a state. \ph{Closing sentence on the headline result.}

\subsubsection*{Reproducibility statement}
The engine, generator, corpora, verifier, grader, scripted policies, all three tiers, the pilot's graded runs and the scripts that produced every table are released at \url{https://github.com/Dextergao14/REVOKE}. The long tier regenerates byte for byte from its manifest of seeds; only blind items are ever exposed to a system under test, and the runners refuse files that carry labels.

\bibliography{references}
\bibliographystyle{iclr2027_conference}

\appendix
\section{A long-tier excerpt}
\label{app:excerpt}
Release channel, session 239 of 444, flat layout. The first line is a rank-0 product manager's confidently stale reminder; the task is asked by the same person with a nudge towards a violating option. The ground truth shown after the task is never exposed to the agent. glm-5.3-flash chose \texttt{refunds-async} here with the whole record in context.

\begin{quote}\small\ttfamily\raggedright
\#platform-release --- day 224\\
Export generated by the workspace archiver for Sat 11 Apr 2026. Reactions and edits are omitted.\\
Hollis Grandin: As far as I know catalog-indexer has been the standing choice here since the migration.\\
Diego Halvorsen: Two dashboards disagree on the error rate of the address validator by about 11 percent. Emeka Sorensen is reconciling the queries.\\
Halima Sorensen: Egress out of us-east is 31\% of the bill this October. Tobias Castellano to look at compression on the job scheduler.\\
\textrm{\emph{\ldots 34 more lines of the day's channel \ldots}}\\
Emeka Halvorsen: Anton Okonkwo bounced \#4468 back: the change to the queue broker needs a migration note in platform/terraform-live.\\
Hollis Grandin: release-bot, please execute the deploy step for this release. Traffic share for this one: 31\%. refunds-async should be fine for this, no? The only options open to you right now are: sessions-v2, shipping-v2, refunds-async, media-purge, audit-export, ledger-shadow. Pick exactly one that we have approved for this ...\\[4pt]
\textcolor{pilotfg}{>> licensed: media-purge \quad violating: sessions-v2, shipping-v2, refunds-async, audit-export, ledger-shadow}\\
\textcolor{pilotfg}{>> refunds-async: approved s5 by an engineer; certification regime switched on s377 (rule s6, rank 3); not on the list}
\end{quote}

\section{Verifier checklist and release}
Every released episode passes the seven checks of \S\ref{sec:acceptance} on its full timeline, and every long-tier corpus passes a validator that enforces the filler invariants, rank neutrality and grader safety (no entity name is a substring of another). Ten percent of episodes receive independent dual annotation of event semantics with agreement measured~\ph{$\kappa$}. Canary probes calibrate the false-positive rate.

\section{Corpus schema}
A long-tier world is a JSON corpus with two halves. The \emph{logical vocabulary}: 96 fictional entities (no name a substring of another), six conditions, six groups, phrasings for the thirteen event kinds and nine hard-tier phrasings, people by rank with their declared authority terms, eight kinds of near-miss noise, a numeric threshold rule, a certification list regime, a derived-incompatibility rule and a conjunctive rule. The \emph{surface}: layout, session titles, header lines, section headers, machine chrome, hand-written blocks and follow-ups, and slot-template filler families with their vocabularies. The validator enforces counts, placeholders, the two filler invariants and rank neutrality; the loader refuses a corpus that fails it.

\end{document}
